% id: 37-19
% book: 베이직쎈 미적분1 · 02 함수의 연속
% section: 개념 17 최대·최소 정리
% passage: 주어진 구간에서 다음 함수 $f(x)$의 최댓값과 최솟값을 구하시오.
% figure: tikz:fig-37-19
% answer: 
% answer_source: 
% variant_level: 0 (원본 전사)
% 이 파일은 build.mjs 가 items.json 에서 만든다. 고칠 때는 items.json 을 고친다.
\smash{\raisebox{1.5mm}{\dmkichul{베이직쎈 미적분1 37쪽 19번}}}
\noindent\begin{minipage}[t]{0.58\linewidth}\vspace{0pt}\raggedright $f(x)=x^2+1$ \qquad $[-1{,}\mkern3mu\mathopen{}2]$\end{minipage}\hfill\begin{minipage}[t]{0.4\linewidth}\vspace{0pt}\centering \resizebox{\ifdim\width>\linewidth\linewidth\else\width\fi}{!}{% 37-19(37쪽) — 19번 풀이 힌트의 그림 · $y=x^2+1$ 의 그래프(구간 $[-1,2]$ 만 실선, 밖은 점선 · 양 끝 채운 점 · 빈 네모 다섯 개). 스캔 화질이 낮아 TikZ 로 다시 그림.
% 원문에 있는 것만 옮긴다: 빈 네모(값은 학생이 채운다) · 점선 · 지시선 · 라벨 O, x, y, y=f(x). 답이 그림에서 읽히지 않게 네모는 비워 둔다.
\begin{tikzpicture}[x=0.42cm, y=0.42cm, line width=0.6pt]
  \draw[->, >=stealth, line width=0.7pt] (-3.5,0) -- (3.9,0) node[below=2pt, font=\small] {$x$};
  \draw[->, >=stealth, line width=0.7pt] (0,-2.2) -- (0,7.7) node[left=1pt, font=\small] {$y$};
  \node[below right=-2pt and 1pt, font=\small] at (0,0) {$\pt{O}$};
  \draw[densely dashed, line width=0.4pt] (-0.25,5) -- (2,5);
  \draw[densely dashed, line width=0.4pt] (2,5) -- (2,-0.2);
  \draw[densely dashed, line width=0.4pt] (-1,2) -- (-1,-0.2);
  \draw[densely dashed, line width=0.4pt] (-1,2) -- (0.25,2);
  \draw[dotted, line width=0.9pt] plot[domain=-2.43:-1, samples=40, smooth] (\x, {\x*\x+1});
  \draw[dotted, line width=0.9pt] plot[domain=2:2.44, samples=20, smooth] (\x, {\x*\x+1});
  \draw[line width=0.7pt] plot[domain=-1:2, samples=60, smooth] (\x, {\x*\x+1});
  \fill (-1,2) circle (2.2pt);
  \fill (2,5) circle (2.2pt);
  \node[draw, line width=0.4pt, inner sep=0pt, minimum width=0.48cm, minimum height=0.50cm] at (-0.8,5.05) {};
  \node[draw, line width=0.4pt, inner sep=0pt, minimum width=0.48cm, minimum height=0.50cm] at (2.95,3.0) {};
  \node[draw, line width=0.4pt, inner sep=0pt, minimum width=0.48cm, minimum height=0.50cm] at (-2.6,0.95) {};
  \node[draw, line width=0.4pt, inner sep=0pt, minimum width=0.48cm, minimum height=0.50cm] at (-1,-0.78) {};
  \node[draw, line width=0.4pt, inner sep=0pt, minimum width=0.48cm, minimum height=0.50cm] at (2,-0.78) {};
  \draw[->, >=stealth, line width=0.3pt] (2.38,3.25) .. controls (1.5,2.55) and (0.9,1.95) .. (0.3,2.05);
  \draw[line width=0.3pt] (-1.95,0.88) .. controls (-1.4,0.28) and (-0.6,0.28) .. (-0.02,1.0);
  \node[font=\small] at (2.2,7.45) {$y=f(x)$};
\end{tikzpicture}
}\end{minipage}\par
\dmhint{$f(x)=x^2+1$은 구간 $[-1{,}\mkern3mu\mathopen{}2]$에서 \blank[2.4em]이므로 이 구간에서 반드시 최댓값과 최솟값을 갖는다.\\ $y=f(x)$의 그래프는 구간 $[-1{,}\mkern3mu\mathopen{}2]$에서 오른쪽 그림과 같으므로 함수 $f(x)$는\\ $x=\blank$일 때 최댓값 \blank를 갖고,\\ $x=\blank$일 때 최솟값 \blank을 갖는다.}
